\chapter{系统设计与实现}

\section{系统总体架构设计}

本系统的整体架构按照职责划分为五个层次，每层承担特定的功能，层与层之间通过明确的接口进行交互，如图\ref{fig:arch}所示。

\begin{figure}[htbp]
    \centering
    \begin{tabular}{|p{12cm}|}
    \hline
    \multicolumn{1}{|c|}{\textbf{Web页面层}：终端识别页面 | 外网用户门户 | 管理员API（FastAPI + Jinja2）} \\
    \hline
    \multicolumn{1}{|c|}{\textbf{识别业务层}：严格识别引擎 | 活体检测服务 | 识别后端适配器} \\
    \hline
    \multicolumn{1}{|c|}{\textbf{摄像头采集层}：采集线程（15FPS） | 识别线程（5FPS） | MJPEG推流} \\
    \hline
    \multicolumn{1}{|c|}{\textbf{数据层}：SQLite数据库 | Repository模式 | 事务管理 | 外键约束} \\
    \hline
    \multicolumn{1}{|c|}{\textbf{配置层}：env.json配置 | 热重载 | 字典式/属性式访问} \\
    \hline
    \end{tabular}
    \caption{系统五层架构设计}
    \label{fig:arch}
\end{figure}

\textbf{配置层}负责管理系统的运行参数。采用基于JSON的配置加载方式，系统启动时读取env.json文件中的配置项，覆盖代码中的默认值。配置项按功能分组，包括Web服务地址和端口、摄像头编号和分辨率、人脸识别阈值、考勤时间窗、活体检测参数等。配置类支持字典式和属性式两种访问方式，并提供配置热重载能力，在运行期间修改env.json后无需重启服务即可生效。

\textbf{数据层}负责SQLite数据库的连接管理、表结构初始化和数据访问封装。系统启动时自动检测数据库文件是否存在，不存在则创建并执行建表脚本。数据访问通过Repository模式封装，上层代码不直接编写SQL语句，而是调用封装好的方法完成增删改查操作。事务管理采用上下文管理器，确保数据写入的原子性。

\textbf{摄像头采集层}负责打开摄像头设备、持续读取视频帧、做镜像和旋转变换，并将帧数据推送给识别线程和MJPEG推流线程。采集线程以15FPS运行，优先保障视频流的流畅度。采集线程不执行识别计算，只负责获取最新的摄像头画面并将其编码为JPEG格式推送给前端。

\textbf{Web页面层}基于FastAPI和Jinja2实现，提供终端识别页面、外网用户门户页面和管理员API接口。终端页面面向站在考勤设备前的使用者，展示实时摄像头画面和识别状态提示。外网用户门户面向远程注册用户，提供注册、登录、人脸采集和审核状态查看功能。

\textbf{识别业务层}是系统的核心，包含严格识别引擎、活体检测服务和识别后端适配器。严格识别引擎接收摄像头帧作为输入，执行一系列门禁判断后返回识别结果。活体检测服务加载OpenVINO模型并执行推理，返回真人概率和屏幕重放风险信号。识别后端适配器统一了YuNet/SFace识别链路的调用接口。

\section{人脸注册流的设计与实现}

\subsection{注册流程设计}

本系统的人脸注册流采用``白名单导入$\to$首次注册$\to$改密$\to$登录$\to$人脸采集$\to$审核$\to$入库''的完整链路。

系统启动时自动读取data/member\_roster.csv文件，将名单中的成员信息同步到roster\_members表中。CSV文件包含姓名、学号/工号、角色、初始密码哈希和状态五个字段。运行期间如果检测到CSV文件发生变化，系统会自动热同步。

用户在名单中的成员首次访问注册页面时，需要输入姓名、学号/工号和初始密码。系统在校验三项信息均与名单匹配后，允许用户设置新密码。注册成功时强制设置符合复杂度要求的新密码，后续登录只使用新密码。注册改密成功后，系统自动跳转到人脸采集页面。

\subsection{动作挑战活体检测}

本系统在人脸注册端采用分步动作挑战机制，替代传统的单张照片上传方式。动作挑战的具体流程如算法\ref{alg:challenge}所示。

\begin{algorithm}[htbp]
    \caption{动作挑战流程}
    \label{alg:challenge}
    \begin{algorithmic}[1]
        \Require 用户浏览器摄像头视频流
        \Ensure 注册照（通过所有校验的人脸图像）
        \State 生成随机3步动作序列 $A = [a_1, a_2, a_3]$ \Comment{如[正脸, 左侧转, 正脸]}
        \State 初始化挑战会话 $S \leftarrow \text{new Session}(A)$
        \For{step $i = 1$ to $3$}
            \State 捕获当前帧图像 $I_i$
            \State 检测人脸并提取landmarks $L_i$
            \State 计算姿态角 $(\text{yaw}, \text{pitch}, \text{roll}) \leftarrow \text{EstimatePose}(L_i)$
            \If{姿态不满足步骤$a_i$的要求}
                \State \Return ``姿态不符合要求''
            \EndIf
            \State 提取人脸编码 $E_i \leftarrow \text{ExtractEncoding}(I_i)$
            \If{$i > 1$ and $\text{CosineDistance}(E_i, E_1) > 0.52$}
                \State \Return ``检测到换人，挑战作废''
            \EndIf
        \EndFor
        \State 从通过挑战的正脸抓拍中按Laplacian清晰度挑选最佳图像 $I_{\text{best}}$
        \State 执行纸张/屏幕翻拍检测 $\text{result} \leftarrow \text{DetectSpoof}(I_{\text{best}})$
        \If{翻拍检测不通过}
            \State \Return ``疑似翻拍，请重新拍摄''
        \EndIf
        \State 保存$ I_{\text{best}} $为pending状态的人脸档案
        \State \Return $I_{\text{best}}$
    \end{algorithmic}
\end{algorithm}

后端从预设模板中随机生成一组3步动作序列，例如``正视镜头$\to$轻微左侧转$\to$正视镜头''。每一步动作对应不同的姿态要求：正脸步骤要求人脸yaw角和roll角均小于阈值，侧转步骤要求yaw角达到一定的偏移量但不超过上限。

用户在每一步完成抓拍后，前端将压缩后的JPEG图像上传到后端。后端对抓拍图执行人脸检测和landmarks提取，计算当前姿态是否满足该步骤的要求。同时提取当前帧的人脸编码，与第一步抓拍的人脸编码进行同人一致性比对。如果检测到换人，当前挑战直接作废。

所有步骤都通过后，后端从通过挑战的正脸抓拍中按Laplacian清晰度评分挑选最佳的一张作为最终注册照。被选中的注册照还需要额外执行纸张和屏幕翻拍检测。最终通过所有校验的注册照保存为pending状态，等待管理员审核。

\subsection{质量校验流水线}

本系统的质量校验采用流水线结构，每个校验器独立负责一种质量检测项。当前实现的校验器包括六类，如表\ref{tab:validators}所示。

\begin{table}[htbp]
    \centering
    \caption{质量校验器设计}
    \label{tab:validators}
    \begin{tabular}{lllc}
        \toprule
        \textbf{校验器} & \textbf{检测指标} & \textbf{阈值} & \textbf{拒绝提示} \\
        \midrule
        单人检测 & 人脸数量 & $=1$ & ``未检测到人脸''/``画面中存在多人'' \\
        人脸尺寸 & 面积占比 & $\geq 5\%$ & ``人脸过小，请靠近摄像头'' \\
        模糊度 & Laplacian方差 & $\geq 100$ & ``图像过于模糊'' \\
        亮度 & 灰度均值 & $[40, 200]$ & ``光线过暗''/``光线过亮'' \\
        姿态 & yaw/pitch/roll & $|\text{angle}| < 30^\circ$ & ``姿态不符合要求'' \\
        翻拍检测 & 边框/纹理/反光 & 多信号命中 & ``疑似翻拍'' \\
        \bottomrule
    \end{tabular}
\end{table}

每个校验器输出统一的结果结构，包含passed（是否通过）、code（错误码）、message（提示信息）和score（评分）。主流程按顺序执行各个校验器，任一校验器不通过则终止流水线并返回失败原因。

纸张/屏幕翻拍检测校验器在最终注册照上执行，检测包裹人脸的矩形边框、疑似屏幕或印刷周期纹理以及低饱和高亮反光区域。该检测器的设计思路是纸质照片和屏幕翻拍往往会在图像中留下可检测的伪影信号，真人直接拍照不会产生这些信号。

\subsection{审核机制}

新上传的人脸档案默认进入pending（待审核）状态，只有管理员审核通过变为approved后，该人脸档案才会被终端识别引擎加载到识别人脸库中。rejected状态的人脸档案不参与识别，用户可以看到驳回原因并重新上传。

审核机制包含以下设计要点。系统内置了十个固定驳回原因选项，包括卡通头像、翻拍屏幕、纸质照片、人脸不清晰、光线异常、姿态错误、遮挡过多、多人脸、信息不符和其他原因。管理员审核时可以多选固定原因，也可以补充自由备注。驳回原因同时保存原因编码和中文标签。

最近三次驳回历史保存在独立的face\_rejection\_history表中。当记录数超过3条时，自动删除最早的记录。保留最近三次驳回历史的目的是让用户知道自己为什么被驳回，有针对性地重新拍照提交。

用户重新上传策略方面，当用户重新提交人脸照片时，系统先完成新照片的质量校验和编码提取，然后删除该用户之前保留的人脸照片文件和旧的face\_profiles记录。系统只为每个用户保留一份当前人脸档案。重新上传存在冷却时间，当前测试值设置为300秒。

\section{实时识别考勤流的设计与实现}

\subsection{严格识别引擎}

本系统的终端识别设计了多道门禁的严格识别引擎。只有所有门禁都通过，才会返回attendance\_ready状态。识别引擎的门禁条件如表\ref{tab:gates}所示。

\begin{table}[htbp]
    \centering
    \caption{严格识别引擎门禁条件}
    \label{tab:gates}
    \begin{tabular}{lll}
        \toprule
        \textbf{门禁} & \textbf{判断条件} & \textbf{不通过返回状态} \\
        \midrule
        单人门禁 & 检测到恰好1张人脸 & no\_face / multi\_face \\
        人脸面积门禁 & 人脸面积占比$\geq$最小门槛 & face\_too\_small \\
        活体前置判断 & 三区间策略通过 & spoof\_detected \\
        编码匹配门禁 & 最小距离$< 0.52$，无歧义 & unknown / ambiguous\_match \\
        多帧稳定确认 & 连续命中$\geq 1$次 & 继续等待 \\
        考勤时间窗 & 当前时间在考勤时段内 & outside\_attendance\_window \\
        冷却时间 & 距离上次签到$\geq 60$秒 & attendance\_cooldown \\
        \bottomrule
    \end{tabular}
\end{table}

单人门禁是识别流程的第一道约束。系统对当前帧执行人脸检测，如果没有检测到人脸，返回no\_face状态。如果检测到多张人脸，返回multi\_face状态。单人入镜的约束确保了识别结果不会产生歧义。

人脸面积门禁在检测到单张人脸后执行。计算人脸区域在整张图像中的面积占比，低于最低门槛时返回face\_too\_small状态。这个门禁的意义在于，当用户距离摄像头过远时，虽然能检测到人脸，但分辨率不足以保证特征提取的准确性。

编码匹配门禁对合格的人脸提取特征编码，与数据库中所有已审核通过的人脸编码逐个计算距离。当最小距离超过阈值（本系统设置为0.52）时，判定为未知用户。如果第一名和第二名候选者的距离比值超过0.85，判定为歧义匹配，避免将用户误认为另一个相似面孔的人。

多帧稳定确认门禁维护一个连续命中计数器，只有同一身份连续命中达到设定次数后，才进入下一阶段。考勤时间窗门禁判断当前时间是否在配置的考勤时间段内。冷却时间门禁判断该用户是否在最近的签到冷却期内，当前配置为60秒。

\subsection{采集与识别线程解耦}

本系统对摄像头采集和识别推理进行了线程级别的解耦设计。旧实现中，摄像头线程在一个循环内同时执行读帧、识别推理、叠字绘制、JPEG编码和MJPEG推流。单次推理耗时在树莓派上可达数百毫秒，识别计算直接阻塞了视频流的刷新，导致终端画面出现明显的卡顿感。

新实现将采集和识别拆分为两个独立的后台线程。采集线程以15FPS持续从摄像头读取最新帧，编码为JPEG后推送给前端MJPEG流。采集线程不做任何识别计算，只负责保持视频流的流畅输出。同时，采集线程会读取识别线程的最新结果并将其叠加到当前帧上，让用户看到实时的识别状态提示。

识别线程以5FPS运行，每次从共享缓冲区获取最新的摄像头帧进行识别推理。识别线程不会追赶旧帧，如果缓冲区中有新的帧到来，识别线程会丢弃尚未处理的旧帧，直接处理最新的一帧。两个线程之间通过线程安全的共享变量交换数据，共享变量的读写使用锁保护。

识别线程还定期（默认每60秒）执行人脸库热重载。热重载期间，识别线程从数据库重新加载所有approved状态的人脸档案，更新内存中的编码库。热重载期间不影响识别线程的正常工作，新的编码库在加载完成后通过原子替换的方式生效。

\section{用户门户与管理接口}

\subsection{用户注册与人脸上传页面}

外网用户门户提供了一组面向普通用户的Web页面，与树莓派终端识别页面彻底分离。门户页面通过访问来源判断入口策略，非本地访问时只保留门户页面入口，本地访问时只保留终端页面入口。

注册页面接收用户输入的姓名、学号/工号和初始密码，提交到后端进行白名单校验。校验通过后允许用户设置新密码，注册成功后自动跳转到人脸采集页面。登录页面接收用户名和新密码，校验通过后创建Session并跳转到用户中心。

用户中心页面展示当前用户的登录状态、人脸档案状态和审核信息。如果用户处于rejected状态，显示驳回原因和重新上传按钮（冷却时间结束后可用）。人脸采集页面提供浏览器摄像头现场拍摄功能，照片经过前端压缩后提交到后端。

\subsection{设备侧管理员API设计}

设备侧管理员API通过FastAPI的APIRouter实现，挂载在\texttt{/api/admin}路径下。所有管理员接口均需要附带Authorization: Bearer Token进行身份验证。API涵盖以下功能模块：设备信息、设备健康、设备性能、成员管理、用户管理、人脸档案和考勤记录。所有API返回数据采用统一的JSON格式。

\subsection{Windows远程管理端对接}

Windows端管理员系统作为一个独立项目运行在已加入Tailscale网络的Windows电脑上。Windows端通过Tailscale内网地址访问树莓派的管理员API，不需要直接暴露树莓派到公网。两个系统通过Tailscale网络实现低耦合连接，树莓派持续独立运行，Windows端可以随时连接或断开。

Windows端当前实现了仪表盘、成员管理、人脸档案审核和考勤记录四个页面。仪表盘展示树莓派健康状态、性能指标和今日考勤汇总。人脸档案审核展示已上传的人脸档案，支持按审核状态筛选和提交审核结果。

\section{数据库设计}

本系统使用SQLite作为本地数据库，核心表结构如表\ref{tab:db}所示。

\begin{table}[htbp]
    \centering
    \caption{数据库表结构}
    \label{tab:db}
    \begin{tabular}{p{2.5cm}p{6cm}p{2cm}}
        \toprule
        \textbf{表名} & \textbf{功能描述} & \textbf{核心字段} \\
        \midrule
        roster\_members & 允许注册的成员名单 & id, name, code, role, status \\
        users & 已登记用户身份信息 & id, name, code, password\_hash \\
        face\_profiles & 人脸特征与审核状态 & id, user\_id, image\_path, encoding, review\_status \\
        attendance\_records & 考勤签到流水记录 & id, user\_id, check\_time, confidence \\
        face\_rejection\_history & 最近三次驳回历史 & id, user\_id, reason\_codes, comment \\
        \bottomrule
    \end{tabular}
\end{table}

roster\_members表作为首次注册的白名单来源，只有名单中的成员才能完成首次注册。face\_profiles表将用户和人脸特征对应起来，终端识别引擎只加载review\_status为approved的记录。attendance\_records表存放实际签到流水，包含签到时间、用户ID、置信度和快照路径。

时间约定方面，所有时间字段统一按设备本地时区写入。实现方式为在Python侧获取当前本地时间后写入数据库，而不是依赖SQLite的CURRENT\_TIMESTAMP默认值，避免UTC时间导致的8小时时差问题。

\section{本章小结}

本章对系统的总体架构和各功能模块进行了详细的设计与实现说明。对五层架构的设计思路和各层职责进行了阐述，说明了模块间解耦的实现方式。对人脸注册流的完整链路、动作挑战活体检测机制、质量校验流水线设计和审核机制进行了详细讲解。对实时识别考勤流的严格识别引擎、采集识别线程解耦进行了说明。对用户门户、管理员API和Windows端对接进行了概述。对数据库表结构和时间约定进行了说明。
